\subsubsection{新39校准与48评价：当前指标通过了哪些验收}
原39校准的982009运行有三项竞争拟合达到200次上限而未收敛，原阈值仍为无穷、评价仍为0/48；追加迭代仅为独立计算诊断。新规则将必要失败后的“保留全部候选”准确编码为零拒绝分数，实际损失、失败和缺失裕量保持原样。规则冻结后改用新的39校准及48评价，旧数据只作开发诊断。

令 $J_a$ 是完整 H+N 的白化SSE，$\kappa\ge0$ 为统一阈值，保留集合为 $\{a:J_a-\min_bJ_b\le\kappa\}$。覆盖成立且拟合完整时，式\eqref{evf:actual-margin}给出
\begin{equation}
 \boxed{s=(-M)_+,\qquad
 M>\kappa:\text{唯一正确};\quad M<-\kappa:\text{误拒};\quad
 |M|\le\kappa:\text{安全未决}.}
 \label{evf:calibrated-margin}
\end{equation}
证明分两种情况：$M\ge0$ 时真候选取得最小损失，最近竞争者差值恰为 $M$；$M<0$ 时最小损失由竞争者取得，真候选的差值为 $-M$。等号保留，得上述三类。必要计算失败另归计算失败，$M$缺失，拒绝分数零只编码未实施拒绝。

新校准取每个独立历史九方法真候选损失差的最大值 $S_i$，再取 $\kappa=\max_{i\le39}S_i=6.034651101192644$。冻结iid混合及覆盖条件下，设下一历史至少一法误拒率为 $R(\kappa)$，则
\begin{equation}
 \Pr_{\rm cal}\{R(\kappa)>.1\}\le.9^{39}=.016423203268\ldots<.05.
 \label{evf:tolerance}
\end{equation}
若风险超过 $.1$，39个独立分数必须全部落在概率质量至多 $.9$ 的下段，独立性即给出该界；并列分数仍按等号保留。此结论针对冻结三末端混合和随机校准集；现场、各子组、多轮及候选覆盖须另验。源和线路条件、种子、字节快照、失败和完整分母均保存。

\begin{table}[htbp]
 \centering\small
 \caption{新独立48历史的完整结果。九方法共用新校准阈值，失败单列。}
 \label{evf:formal48}
 \begin{tabular}{lrrrr}
 \toprule
 方法 & 唯一正确 & 安全未决 & 误拒 & 计算失败\\
 \midrule
同长度无EV & 9 & 39 & 0 & 0\\
固定点恒充 & 28 & 20 & 0 & 0\\
简单计量质量曲线 & 26 & 22 & 0 & 0\\
固定点信息曲线 & 27 & 21 & 0 & 0\\
随机点信息曲线 & 29 & 19 & 0 & 0\\
移动信息曲线 & 30 & 18 & 0 & 0\\
固定点V2G & 27 & 19 & 2 & 0\\
移动V2G & 27 & 20 & 1 & 0\\
人员扩窗单向 & 28 & 20 & 0 & 0\\
\bottomrule
 \end{tabular}
\end{table}

固定信息曲线相对恒充平均实际裕量提高 $268.523$，配对SE为 $96.630$（48独立历史）；唯一正确却由28变27，配对1胜2负。移动相对固定提高 $111.017\pm78.233$（均值与SE），唯一正确4胜1负、总计30/48。固定V2G相对单向提高 $270.287\pm124.955$，唯一正确1胜1负且误拒2；人员扩窗相对移动为 $-8.577\pm6.093$，唯一正确0胜2负。平均裕量与跨过筛选阈值的次数是不同终点，当前固定曲线优化未通过主要排错验收。

九方法中至少一法误拒2/48，95\%二项区间 $[.00509,.14254]$。弱支路8历史中，固定/移动单向各仅1次唯一正确；star7历史单向全未决，V2G另有误拒。局部四候选总体内未漏真；更大WZZT候选覆盖仍需检验。

固定点后续价格在22/42条成功设计中提高分数，平均提高1.598；最佳操作优化初值相对恒充平均提高24.454。分数增量与实际裕量增量的描述相关，固定相对恒充为 $.086$、移动相对恒充为 $.028$；分别11/42、12/42条分数增量正而实际裕量增量负。全部48分母中，另有4未启用、2设计回退。随机方法21次选址搜索失败及2次共同回退均保留。

指标缺口有明确原因：当前 $I$ 比较疑似边与零边，并未把另一正cherry的内部边作为竞争补偿。两种内部边导数相同即可对star可见而相互不可区分。985001、985037的star历史解释又在支持规则外，无法建立当前收缩展开点，保留恒充回退。因此下一阶段应在合法正候选上构造成对历史和未来补偿，先核查已有曲线距离与价格，再接入优化；弱边继续未决，不为提高分数缩小竞争域。

新39校准独立重放156历史拟合、1404保存联合拟合和351真实响应；48评价重放192历史拟合、1728保存联合拟合和432真实响应，全部通过。23份冻结源码及37安全规则测试来源字节一致，原始与紧凑记录科学结果相同。算法AC调用分别89379和109103，排除真值生成。九方法带缓存benchmark均值分别22.569和28.020秒；选择超过15秒为9/39、11/48，整历史超过90秒为0。该时间是完整对照benchmark，不是单策略部署时间。校准保存计时发现363/1122个跨设备域精确重复价格请求；共享缓存可能减少工作，实际加速另行验证。

本批连续功率假定为0--6 kW，人员为外生准备接口窗口、接线时长0、返程耗电0。实际设备最小开机功率、追踪、非零接线及返程需进入下一操作域。下一阶段先核查正边竞争、信息与排错的联系及精确缓存，再冻结新策略和新独立数据；滚动、人员路线、扩网、proper预测EIG与全局AC/EV最优证书保持未完成。已有负例与原校准失败全部保留。
